% P8 capability-gap taxonomy: the four LLM wins (task categories, not accuracy deltas)
\section{The Capability-Gap Taxonomy: Where the LLM Wins}
\label{sec:p8-taxonomy}

The LLM's four wins share a structure: each is a \emph{task category the tree
cannot enter}, not an accuracy improvement on a task the tree already
performs. A tree cannot be beaten at reading a forged utility bill, drafting a
SAR narrative, reasoning about an unlabeled typology, or interviewing an alert
queue---because it cannot attempt any of them. We take the four in turn.

\subsection{Gap 1: Document and Image Fraud (Perception)}
\label{sec:p8-gap-docs}

A persistent channel of fraud runs through unstructured evidence:
doctored receipts submitted in dispute claims, tampered identity documents at
onboarding, synthetic or altered checks, screenshots fabricated for
first-party fraud claims. These artifacts are raw pixels and embedded
text---inputs that lie outside any tree's feature space by construction. The
tree sees such evidence only through whatever handcrafted features an upstream
pipeline extracts, and fraudsters attack precisely the residual the pipeline
discards. Vision--language models change the input space itself: multimodal
LLMs can detect and localize image forgeries while producing textual
justifications \citep{xu2024fakeshield}, and systematic benchmarks now
evaluate large vision--language models across more than a hundred forgery
types \citep{wang2025forensicsbench}. The benchmark literature is equally
clear that current VLMs are far from solved on comprehensive forgery
detection \citep{wang2025forensicsbench}---our claim is not that the VLM is
accurate enough to auto-decide, but that it is the \emph{only} model family in
the pipeline that can look at the evidence at all. In the architecture of
\secref{sec:p8-architecture} its output is a feature vector and a written
observation, not a verdict.

\subsection{Gap 2: Compliance Narration (The Scribe)}
\label{sec:p8-gap-narration}

United States banking regulation obliges institutions to produce
\emph{prose}. A Suspicious Activity Report must be filed for qualifying
transactions \citep{cfr1020_320}, and FinCEN's guidance is explicit that the
narrative section must tell the who/what/when/where/why of the suspicion in
clear, complete text \citep{fincen2003sarnarrative}. On the lending side,
ECOA/Regulation~B requires that adverse-action notices state specific,
principal reasons \citep{cfr1002_9}, and the CFPB has affirmed that model
complexity is no defense \citep{cfpb2022circular}. These are text-generation
obligations with legal force, currently discharged by analysts writing under
backlog pressure. A tree emits a float; it cannot enter this category at any
accuracy. An LLM conditioned on the case file, the tree's feature attributions,
and the regulatory template can draft both document types for human review,
and early agentic systems now target SAR-narrative drafting specifically, with
human-in-the-loop review \citep{naik2025coinvestigator}.
\secref{sec:p8-compliance} grounds this gap in the primary sources, because it
is the paper's most durable claim: it survives any future accuracy result.

\subsection{Gap 3: Cold-Start on Novel Typologies (Few-Shot)}
\label{sec:p8-gap-coldstart}

A supervised tree is, by definition, a summary of labeled history. When a new
fraud typology emerges---a novel scam pattern, a new mule-recruitment
mechanic, an exploit of a just-launched product feature---the tree has zero
training signal until enough confirmed cases accumulate, which at fraud base
rates and confirmation latencies \citep{dalpozzolo2018realistic} means weeks
of unpriced exposure. An LLM can operate in exactly this window: given a
typology description from a fraud-intelligence bulletin and a handful of
exemplar cases, it can flag candidate matches few-shot, before any labeled
dataset exists---the regime where LLM-over-tabular approaches are genuinely
competitive \citep{hegselmann2023tabllm}, and where in-context reasoning over
serialized financial evidence has shown analyst-style red-flag detection
without task-specific training \citep{pirmorad2025icl_aml}. The LLM here is
not a better classifier; it is a stopgap classifier \emph{that exists} during
the label vacuum, whose flagged cases then become the seed labels that bring
the tree up to speed.

\subsection{Gap 4: Agentic Alert-Queue Investigation}
\label{sec:p8-gap-agentic}

Between the scorer's alert and the analyst's disposition sits hours of
mechanical investigation: pulling account history, cross-referencing devices
and merchants, checking prior disputes, summarizing the case. An LLM agent can
execute this gather-and-summarize loop over the alert queue---the pattern
emerging in agentic AML investigation systems \citep{naik2025coinvestigator,
pirmorad2025icl_aml}. In our internal estimate, an agentic pass over one alert
costs roughly $85\times$ less than the loaded cost of the equivalent
human-analyst time; we stress that this figure is an \emph{internal, assumption-laden
estimate} (token prices, analyst compensation, and per-alert
investigation time as of our June~2026 records), reported to one significant
figure and not validated against any production deployment. Even at an order
of magnitude worse, the economics reshape the queue: the agent does not
replace the analyst but re-sorts and pre-digests the queue so that human
minutes land on the cases where they change the outcome. Crucially this role
is \emph{post-score}: the agent operates on alerts the tree has already
raised, outside the real-time path and behind the injection-hardening that
\secref{sec:p8-scorer} argues the synchronous path must never expose.

\paragraph{Summary.} The taxonomy's lesson is that ``LLM vs.\ XGBoost'' was a
category error. The four wins are perception, narration, few-shot
generalization, and agency---four capabilities bundled under the label ``AI''
that the scoring benchmark never measured. The next section assembles the
division of labor that the taxonomy implies.
